% Ad Atticum 5.20 (ad-atticum-05-20) --- English translation
% Translated by Claude (Anthropic) from the Latin (Perseus).
% Style guide: STYLE.md.

\ciceroLetterOpener{CICERO ATTICO salutem}

\ciceroSection{1} On the morning of the Saturnalia the people of Pindenissus surrendered themselves to me, on the fifty-seventh day from the day on which we began to besiege them. ``The devil!'' you will say --- ``who in the world are these Pindenissites? I have never heard the name.'' What am I to do? Could I turn Cilicia into Aetolia or Macedonia? Take it now from me that with this army --- and in this country --- so great a piece of work could not have been brought off. Here it is \textit{[Greek: in summary]}: which is what you grant me in your most recent letter. How I arrived at Ephesus you know --- you sent me your own congratulations on the splendour of that day, than which nothing has ever pleased me more. From there, in the towns through which we passed, having been received with wonderful welcome, we came to Laodicea on the day before the Kalends of Sextilis. There we delayed two days, were very much in the limelight, and by the honour of our reception undid all the earlier wrongs of the country; the same again at Colossae; then we stayed five days at Apamea, three at Synnada, five at Philomelium, ten at Iconium. Nothing more even than that jurisdiction, nothing milder, nothing more grave.

\ciceroSection{2} Then I came into camp on the seventh day before the Kalends of September. On the third day I held the muster of the army at Iconium. From this camp, when grave reports of the Parthians began to arrive, I pushed on into Cilicia by way of that part of Cappadocia which touches Cilicia, with this design: that Armenian Artavasdes, and the Parthians themselves, should reckon themselves shut off from Cappadocia. When I had kept the camp for five days at Cybistra of Cappadocia, I was informed for certain that the Parthians were far off from that approach into Cappadocia, and that they were threatening Cilicia more nearly. So at once I made the march into Cilicia by the Gates of the Taurus. I came to Tarsus on the third day before the Nones of October.

\ciceroSection{3} From there I hurried on to the Amanus, which divides Syria from Cilicia along the watershed: a mountain country full of our eternal enemies. There, on the third day before the Ides of October, we cut down a great number of the enemy. Strongholds, sound as they were, we took --- Pomptinus coming up at night, ourselves at morning --- and fired them. We were hailed as imperators. We held camp for a few days in the very place where Alexander had pitched against Darius at Issus --- an imperator no little better than either you or I. Having stayed there five days, after the Amanus had been overrun and laid waste, we set off. In the meantime --- for you know how there is talk of \textit{[Greek: panic-fears]}, and likewise of \textit{[Greek: the empty alarms of war]} --- on the rumour of our coming, Cassius (who was being held shut up in Antioch) took courage, and the Parthians took fright. So as they were withdrawing from the town Cassius pursued them and managed it well. In that flight the Parthian general Osaces took a wound with great authority on his side, and died of it a few days afterwards. Our name stood high in Syria.

\ciceroSection{4} In the meantime Bibulus arrived; he was eager, I take it, to be on equal terms with us by means of this empty title. In the same Amanus he began to look for his little laurel in a baker's tart. But he lost the whole first cohort, and the chief centurion --- a man of his own rank, Asinius Dento --- and the rest of that cohort, and Sex. Lucilius, the son of T. Gavius Caepio, a rich and splendid man, a tribune of the soldiers. He had taken a heavy blow indeed, both for what it was and for the moment of it. We pressed on to Pindenissus, which was the most strongly walled town of the Eleutherocilicians in all memory, and was in arms. Wild men, fierce, prepared in every respect for self-defence. We hemmed them in with rampart and ditch; with a huge mound, mantlets, a most lofty tower, a great supply of artillery, many archers, with great labour and great apparatus, with many of our own wounded but the army whole, we finished the business. A merry Saturnalia indeed for the men --- to whom (the horses excepted) we made over the rest of the plunder. The slaves were being sold off on the third day of Saturnalia. As I write this, the bidding at the tribunal stands at HS 120,000. From here I am handing over to my brother Quintus the army, to be led off to winter quarters in an ill-pacified country; I myself am withdrawing to Laodicea.

\ciceroSection{6} So much for what is current. But let me come back to what is past. As to what you most strenuously urge upon me, and what is worth more than all else --- that on which you take such pains, that we may give satisfaction even to your Ligurian \textit{[Greek: Momus]} \textit{[fault-finder]} --- I will die if anything could be done more elegantly. I do not even any more call it \textit{continentia}, that virtue which seems to resist mere pleasure. In my whole life I have never been affected by any pleasure so great as I am by this integrity; and it is not the renown --- which is great --- that delights me, so much as the thing in itself. What more would you have? It was worth it. I did not know myself, and did not know well enough what I could do in this line. \textit{[Greek: I am rightly puffed up]}. Nothing is dearer to me than this. Meanwhile here are the \textit{[Greek: bright doings]}. Ariobarzanes lives and reigns by my labours: \textit{[Greek: in]} counsel and authority, and because I made myself \textit{[Greek: inapproachable]} to his plotters, not merely \textit{[Greek: incorruptible]} ---- I have saved the king and the kingdom. Meanwhile from Cappadocia not so much as a single hair. Brutus, who was downcast, I have raised up as much as I could --- a man I love no less than you do, almost as much, I was going to say, as I love you. And further, I hope that for the whole year of our command there will not be a single farthing of expense to the province.

\ciceroSection{7} There you have everything. Now I am preparing to send a public despatch to Rome. It will be fuller than if I had sent it from the Amanus. To think that you will not be at Rome! But the whole thing turns on what is to be on the Kalends of March. I am afraid lest, when the matter of the provinces comes up, if Caesar resists, I shall be kept on. If you were there, I should fear nothing.

\ciceroSection{8} I come back to city news, which, long unaware of it, I have at last gathered from your most welcome letters, on the fifth day before the Kalends of January. Your freedman Philogenes took every care to have them brought through, by a very long and not very safe road. Those you say you gave to Laenius's slaves I had not received. Pleasant the news about Caesar, both what the Senate decreed and what you hope for. If he gives way to all this, we are safe. That Seius was singed in the Plaetorian conflagration I take less ill. Why Lucceius was so vehement against Q. Cassius, and what was done, I am eager to know.

\ciceroSection{9} When I have come to Laodicea, I am bidden to give the plain toga to Quintus, your sister's son; whom I shall school the more carefully. Deiotarus --- whose great auxiliaries I made much use of --- was, as he wrote, on the point of coming to me with the boys to Laodicea. I await also your Epirote letter, that I may have a reckoning not only of your business but also of your leisure. Nicanor is on duty, and is liberally treated by me. I think I shall send him to Rome with the public despatch, both so that it may be the more diligently delivered, and so that he may report back to me with certain news of you, and from you. That Alexis so often adds the salutation to me is gratifying; but why does he not do in his own letter what my Alexis does in writing to you? For Phemio a \textit{[Greek: horn]} \textit{[lyre]} is being sought. But enough of this. Take care of yourself, and let me know when you are thinking of Rome. Again and again, farewell.

\ciceroSection{10} Your people and your interests I had commended most carefully to Thermus, both in person at Ephesus and now by letter; and I have understood that the man himself is most devoted to you. As for what I wrote you before about the house of Pammenes, please see to it, by every means, that what the boy holds by your favour and mine is not in any way wrested from him. I think it honourable for both of us; and to me it will be most welcome.
